% GENERATED FILE. Edit canonical lessons, then run npm run generate:books.
% canonical-source-hash: fnv1a64:f51b14eb89e961e7
% canonical-lessons: LA-C436-puppis, LA-C436-prora, LA-C436-carina, LA-C436-transtrum, LA-C436-onerarius

\chapter{Stern, Prow, Keel, Rowing bench, For cargo}
\label{ch:la-stern-prow-keel-rowing-bench-for-cargo}
\hlchaptermodality{436}

\begin{chapteropening}
\textbf{By the end of this chapter:} \emph{I can name the parts of a ship: the stern, the prow, the keel, the rowing benches, and say a ship is for cargo.}

Complete the last lesson of chapter 436: I can name the parts of a ship: the stern, the prow, the keel, the rowing benches, and say a ship is for cargo.
\end{chapteropening}

\section[puppis, puppis]{puppis, puppis — the stern (of a ship); a ship}

\label{lesson:LA-C436-puppis}



\begin{quote}
Before the new one: say the Latin for a punt, then the Latin for a raft.
\end{quote}

\subsection*{You'll want to know: puppis, puppis}
\textbf{puppis, puppis} — \textquotedblleft{}the stern (of a ship); a ship\textquotedblright{}.

\begin{quote}
Gubernātor in puppī stat et nāvem regit.
\end{quote}

\begin{quote}
\emph{The helmsman stands at the stern and steers the ship.}
\end{quote}

Third declension, feminine, with accusative \emph{puppim} and ablative \emph{puppī}. English \emph{poop deck} comes from it through French.

\begin{cousinweb}
The origin of \emph{puppis} is not settled: Lewis \& Short call it doubtful and EtymDB gives no parent. English \emph{poop}, the raised deck at the stern, comes from it through French \emph{poupe}.
\end{cousinweb}

\begin{grammarlens}[title={What you've built}]
One more word for telling, explaining and describing at B1.
\end{grammarlens}

\subsection*{Your turn}
\begin{itemize}
\raggedright
  \item \emph{Say it:} \emph{puppis, puppis}
  \item \emph{Say it:} \emph{puppis, puppis}, once more
  \item \emph{Say it:} a sentence of your own with \emph{puppis, puppis}
  \item \emph{Recall:} say the Latin for a punt, then the Latin for a raft, then say \emph{puppis, puppis} again
\end{itemize}

\begin{tcolorbox}[breakable,colback=teal!4,colframe=teal!35!black,title={Before you move on}]
What does puppis, puppis mean? (\textquotedblleft{}The stern (of a ship); a ship\textquotedblright{}.) Say it once more.
\end{tcolorbox}
\section[prōra, prōrae]{prōra, prōrae — the prow, the bow (of a ship)}

\label{lesson:LA-C436-prora}



\begin{quote}
Before the new one: say the Latin for a raft, then the Latin for the stern.
\end{quote}

\subsection*{You'll want to know: prōra, prōrae}
\textbf{prōra, prōrae} — \textquotedblleft{}the prow, the bow (of a ship)\textquotedblright{}.

\begin{quote}
Nauta in prōrā stat et lītus spectat.
\end{quote}

\begin{quote}
\emph{A sailor stands in the prow and watches the shore.}
\end{quote}

First declension, feminine: \emph{prōra, prōrae f.} \emph{Prōra et puppis}, the whole of a thing, front to back. English \emph{prow} comes from it through French.

\begin{cousinweb}
\emph{prōra} was borrowed from Greek \emph{prōira} 'prow' (EtymDB), built on the root of \emph{pro} 'forward'. English \emph{prow} comes from it through French \emph{proue}.
\end{cousinweb}

\begin{grammarlens}[title={What you've built}]
One more word for telling, explaining and describing at B1.
\end{grammarlens}

\subsection*{Your turn}
\begin{itemize}
\raggedright
  \item \emph{Say it:} \emph{prōra, prōrae}
  \item \emph{Say it:} \emph{prōra, prōrae}, once more
  \item \emph{Say it:} a sentence of your own with \emph{prōra, prōrae}
  \item \emph{Recall:} say the Latin for a raft, then the Latin for the stern, then say \emph{prōra, prōrae} again
\end{itemize}

\begin{tcolorbox}[breakable,colback=teal!4,colframe=teal!35!black,title={Before you move on}]
What does prōra, prōrae mean? (\textquotedblleft{}The prow, the bow (of a ship)\textquotedblright{}.) Say it once more.
\end{tcolorbox}
\section[carīna, carīnae]{carīna, carīnae — the keel (of a ship); a boat}

\label{lesson:LA-C436-carina}



\begin{quote}
Before the new one: say the Latin for the stern, then the Latin for the prow.
\end{quote}

\subsection*{You'll want to know: carīna, carīnae}
\textbf{carīna, carīnae} — \textquotedblleft{}the keel (of a ship); a boat\textquotedblright{}.

\begin{quote}
Nāvis nova carīnam longam et validam habet.
\end{quote}

\begin{quote}
\emph{The new ship has a long, strong keel.}
\end{quote}

First declension, feminine: \emph{carīna, carīnae f.} First a nutshell, then the keel. English \emph{careen} 'to tip a ship on its side' comes from it through French.

\begin{cousinweb}
\emph{carīna} is inherited and is usually traced to PIE \emph{\'{k}erh\textsubscript{2}-} (EtymDB); Lewis \& Short compare Greek \emph{karyon} 'nut', since its first sense is a nutshell. English \emph{careen} comes from it through French \emph{carène}.
\end{cousinweb}

\begin{grammarlens}[title={What you've built}]
One more word for telling, explaining and describing at B1.
\end{grammarlens}

\subsection*{Your turn}
\begin{itemize}
\raggedright
  \item \emph{Say it:} \emph{carīna, carīnae}
  \item \emph{Say it:} \emph{carīna, carīnae}, once more
  \item \emph{Say it:} a sentence of your own with \emph{carīna, carīnae}
  \item \emph{Recall:} say the Latin for the stern, then the Latin for the prow, then say \emph{carīna, carīnae} again
\end{itemize}

\begin{tcolorbox}[breakable,colback=teal!4,colframe=teal!35!black,title={Before you move on}]
What does carīna, carīnae mean? (\textquotedblleft{}The keel (of a ship); a boat\textquotedblright{}.) Say it once more.
\end{tcolorbox}
\section[trānstrum, trānstrī]{trānstrum, trānstrī — a rowing bench, a thwart; a crossbeam}

\label{lesson:LA-C436-transtrum}



\begin{quote}
Before the new one: say the Latin for the prow, then the Latin for the keel.
\end{quote}

\subsection*{You'll want to know: trānstrum, trānstrī}
\textbf{trānstrum, trānstrī} — \textquotedblleft{}a rowing bench, a thwart; a crossbeam\textquotedblright{}.

\begin{quote}
Rēmigēs in trānstrīs sedent.
\end{quote}

\begin{quote}
\emph{The rowers sit on the benches.}
\end{quote}

Second declension, neuter: \emph{trānstrum, trānstrī n.} The bench runs across the ship. English \emph{transom} is usually traced to it.

\begin{cousinweb}
\emph{trānstrum} is built on \emph{trāns} 'across' with the suffix \emph{-trum} (EtymDB; Lewis \& Short): a beam laid across a ship. English \emph{transom} is usually traced to it.
\end{cousinweb}

\begin{grammarlens}[title={What you've built}]
One more word for telling, explaining and describing at B1.
\end{grammarlens}

\subsection*{Your turn}
\begin{itemize}
\raggedright
  \item \emph{Say it:} \emph{trānstrum, trānstrī}
  \item \emph{Say it:} \emph{trānstrum, trānstrī}, once more
  \item \emph{Say it:} a sentence of your own with \emph{trānstrum, trānstrī}
  \item \emph{Recall:} say the Latin for the prow, then the Latin for the keel, then say \emph{trānstrum, trānstrī} again
\end{itemize}

\begin{tcolorbox}[breakable,colback=teal!4,colframe=teal!35!black,title={Before you move on}]
What does trānstrum, trānstrī mean? (\textquotedblleft{}A rowing bench, a thwart; a crossbeam\textquotedblright{}.) Say it once more.
\end{tcolorbox}
\section[onerārius, onerāria, onerārium]{onerārius, onerāria, onerārium — for cargo, for carrying loads}

\label{lesson:LA-C436-onerarius}



\begin{quote}
Before the new one: say the Latin for the keel, then the Latin for a rowing bench.
\end{quote}

\subsection*{You'll want to know: onerārius, onerāria, onerārium}
\textbf{onerārius, onerāria, onerārium} — \textquotedblleft{}for cargo, for carrying loads\textquotedblright{}.

\begin{quote}
Nāvēs onerāriae frūmentum ad urbem vehunt.
\end{quote}

\begin{quote}
\emph{Cargo ships carry grain to the city.}
\end{quote}

An adjective of three endings. \emph{Nāvis onerāria}, a merchant ship, or simply \emph{onerāria}; \emph{iūmentum onerārium}, a pack animal.

\begin{cousinweb}
\emph{onerārius} is built on \emph{onus, oneris} 'a load' with the suffix \emph{-ārius} (EtymDB; Lewis \& Short). English \emph{onerary} 'fit for loads' is a rare borrowing of it.
\end{cousinweb}

\begin{grammarlens}[title={What you've built}]
One more word for telling, explaining and describing at B1.
\end{grammarlens}

\subsection*{Your turn}
\begin{itemize}
\raggedright
  \item \emph{Say it:} \emph{onerārius, onerāria, onerārium}
  \item \emph{Say it:} \emph{onerārius, onerāria, onerārium}, once more
  \item \emph{Say it:} a sentence of your own with \emph{onerārius, onerāria, onerārium}
  \item \emph{Recall:} say the Latin for the keel, then the Latin for a rowing bench, then say \emph{onerārius, onerāria, onerārium} again
\end{itemize}

\begin{tcolorbox}[breakable,colback=teal!4,colframe=teal!35!black,title={Before you move on}]
What does onerārius, onerāria, onerārium mean? (\textquotedblleft{}For cargo, for carrying loads\textquotedblright{}.) Say it once more.
\end{tcolorbox}
